\chapter*{DANH MỤC THUẬT NGỮ}
\addchapternopage{DANH MỤC THUẬT NGỮ}

\begin{longtable}{|l|p{9.5cm}|}
    \hline
    \textbf{Thuật ngữ tiếng Anh} & \textbf{Thuật ngữ tiếng Việt tương đương / Định nghĩa} \\
    \hline
    \endfirsthead
    \hline
    \textbf{Thuật ngữ tiếng Anh} & \textbf{Thuật ngữ tiếng Việt tương đương / Định nghĩa} \\
    \hline
    \endhead
    \hline
    \endfoot
    Embedding & Nhúng (Quá trình biểu diễn thực thể thành vector số học liên tục) \\
    \hline
    Vector Database & Cơ sở dữ liệu Vector (Hệ thống lưu trữ và tìm kiếm vector tương đồng) \\
    \hline
    Tech stack & Hệ thống nền tảng công nghệ / Bộ công nghệ triển khai (Tập hợp các thư viện, framework và hạ tầng) \\
    \hline
    Prompt & Lời nhắc (Prompt) (Ngữ cảnh, câu hỏi và chỉ thị đầu vào cho mô hình ngôn ngữ) \\
    \hline
    Hallucination & Hiện tượng ảo giác (Mô hình sinh thông tin sai lệch thiếu căn cứ) \\
    \hline
    Chunking & Phân đoạn / Phân mảnh (Chia nhỏ văn bản thành các phân đoạn độc lập) \\
    \hline
    Information Retrieval & Truy xuất thông tin (IR) \\
    \hline
    Fine-tuning & Tinh chỉnh mô hình (Huấn luyện chuyên sâu trên tập dữ liệu đặc thù) \\
    \hline
    Generative AI & Trí tuệ nhân tạo tạo sinh \\
    \hline
    Hybrid retrieval & Truy xuất lai (Kết hợp đồng thời kênh từ khoá và kênh ngữ nghĩa) \\
    \hline
    Dense / Sparse retrieval & Truy xuất ngữ nghĩa theo vector đặc / Truy xuất từ khóa theo vector thưa \\
    \hline
    Re-ranking & Tái xếp hạng (Đánh giá lại điểm số và sắp xếp tập ứng viên sơ tuyển) \\
    \hline
    Cross-encoder & Bộ mã hoá chéo / Mô hình tương tác chéo (Transformer xử lý đồng thời cặp truy vấn -- tài liệu) \\
    \hline
    Bi-encoder & Kiến trúc hai nhánh mã hóa độc lập / Bi-encoder (Mô hình mã hoá truy vấn và tài liệu tách biệt) \\
    \hline
    Relevance gate & Cổng lọc độ liên quan (Cơ chế lọc ứng viên dựa trên ngưỡng khoảng cách ngữ nghĩa) \\
    \hline
    Ground truth & Dữ liệu chuẩn đối sánh / Nhãn chuẩn xác thực \\
    \hline
    Ablation study & Thực nghiệm bóc tách thành phần (Phân tích đóng góp độc lập của từng cấu phần) \\
    \hline
    LLM-as-a-judge & Phương pháp đánh giá chất lượng bằng mô hình ngôn ngữ lớn độc lập \\
    \hline
    Human-in-the-loop & Quy trình phối hợp chuyên gia rà soát (Chuyên gia trực tiếp thẩm định và kiểm định) \\
    \hline
    Layout segmentation & Phân vùng bố cục trang tài liệu \\
    \hline
    Multimodal RAG & Sinh văn bản tăng cường truy xuất đa phương thức (Tích hợp đồng thời học liệu văn bản và hình ảnh) \\
    \hline
    Late interaction & Cơ chế tương tác trễ (ColBERT tính tương đồng ở mức token sau khi mã hóa ngữ cảnh) \\
    \hline
    Zero-shot retrieval & Truy xuất không cần tinh chỉnh mẫu (Khả năng tìm kiếm tổng quát trên tập dữ liệu mới) \\
    \hline
    Hit@K / Recall@K & Tỷ lệ trúng tại K / Độ phủ tại K (Tỷ lệ câu hỏi tìm thấy tài liệu liên quan trong top K kết quả) \\
    \hline
    Out-of-domain & Ngoài phạm vi kho tri thức (Các câu hỏi không thuộc nội dung sách giáo khoa mục tiêu) \\
    \hline
    Formula degradation & Hiện tượng suy biến công thức (Lỗi phân rã ký hiệu toán -- lý -- hóa và chỉ số dưới qua OCR) \\
    \hline
    Line stitching & Ghép dòng lai phân cấp (Quy trình hợp nhất có điều kiện giữa văn bản OCR và công thức chuyên sâu) \\
    \hline
\end{longtable}